\subsection{DIRECT SYSTEMS OF MAPPINGS}

PROPOSITION 6. Let I be a directed set, let \((\mathbf{E}_{\alpha}, f_{\beta \alpha})\) be a direct system of sets relative to I, let \(\mathbf{E} = \lim \mathbf{E}_{\alpha}\) be the direct limit, and for each \(\alpha \in \mathbf{I}\) let

\[
f _ {\alpha}: \mathbf {E} _ {\alpha} \rightarrow \mathbf {E}
\]

be the canonical mapping. For each \(\alpha \in \mathbf{I}\) , let \(u_{\alpha}\) be a mapping of \(\mathbf{E}_{\alpha}\) into a set \(\mathbf{F}\) such that

\[
u _ {\beta} \circ f _ {\beta x} = u _ {\alpha}\tag{23}
\]

whenever \(\alpha \leqslant \beta\) .

Then:

\begin{enumerate}
\def\labelenumi{(\alph{enumi})}
\tightlist
\item
  There exists a unique mapping \(u\) of \(\mathbf{E}\) into \(\mathbf{F}\) such that
\end{enumerate}

\[
u _ {\alpha} = u \circ f _ {\alpha}\tag{24}
\]

\[
{\text{for all}} \alpha \in \mathbf {I}.
\]

\begin{enumerate}
\def\labelenumi{(\alph{enumi})}
\setcounter{enumi}{1}
\item
  \(u\) is surjective if and only if \(\mathbf{F}\) is the union of the sets \(u_{\alpha}(\mathbf{E}_{\alpha})\) .
\item
  \(u\) is injective if and only if for each \(\alpha \in I\) the relations \(x \in E_{\alpha}\) , \(y \in E_{\alpha}\) , \(u_{\alpha}(x) = u_{\alpha}(y)\) imply that there exists \(\beta \geqslant \alpha\) such that \(f_{\beta \alpha}(x) = f_{\beta \alpha}(y)\) .
\item
  With the notation of no. 5, let \(v\) be the mapping of \(G\) into \(F\) which agrees with \(u_{\alpha}\) on \(E_{\alpha}\) for each \(\alpha \in I\) (Chapter II, §4, no. 7, Proposition 8). The hypothesis implies that \(v\) is compatible with the equivalence relation \(R\) (Chapter II, §6, no. 5); hence there exists a unique mapping \(u\) of \(E = G / R\) into \(F\) such that \(v = u \circ f\) (loc. cit.).
\item
  Since \(\mathbf{E}\) is the union of the \(f_{\alpha}(\mathbf{E}_{\alpha})\) , the relation \(\mathbf{F} = \bigcup_{\alpha \in \mathbf{I}} u_{\alpha}(\mathbf{E}_{\alpha})\) is clearly necessary and sufficient for \(u\) to be surjective.
\item
  By Lemma 1 of no. 5 any two elements of E can always be written in the form \(f_{\alpha}(x)\) and \(f_{\alpha}(y)\) , where \(x \in \mathbf{E}_{\alpha}\) and \(y \in \mathbf{E}_{\alpha}\) , for a suitable choice of \(\alpha \in \mathbf{I}\) . It follows also from the lemma that the relation \(f_{\alpha}(x) = f_{\alpha}(y)\) is equivalent to the existence of \(\beta \geqslant \alpha\) such that \(f_{\beta \alpha}(x) = f_{\beta \alpha}(y)\) . Since \(u_{\alpha}(x) = u(f_{\alpha}(x))\) and \(u_{\alpha}(y) = u(f_{\alpha}(y))\) , this completes the proof.
\end{enumerate}

If the mapping \(u\) is bijective, it is sometimes said, by abuse of language, that \(F\) is the direct limit of the family \((\mathbf{E}_{\alpha})\) .

Remark. Suppose that each of the mappings \(f_{\beta \alpha}\) is injective. Then each of the \(f_{\alpha}\) is injective, by the definition of the relation R. In this case we generally identify \(E_{\alpha}\) and \(f_{\alpha}(E_{\alpha})\) and consider E therefore as the union of the \(E_{\alpha}\) . Conversely, let \((F_{\alpha})_{\alpha \in I}\) be an increasing family of subsets of a set F and suppose that F is the union of this family. If \(j_{\beta \alpha}\) denotes the canonical injection of \(F_{\alpha}\) into \(F_{\beta}\) for \(\alpha \leqslant \beta\) , then it follows from Proposition 6 that we may identify F with the direct limit of the family \(F_{\alpha}\) with respect to the family of mappings \((j_{\beta \alpha})\) , and the canonical mapping of \(F_{\alpha}\) into \(\lim F_{\alpha}\) with the canonical injection of \(F_{\alpha}\) into F, for each \(\alpha \in I\) .

COROLLARY 1. Let \((\mathbf{E}_{\alpha}, f_{\beta \alpha})\) and \((\mathbf{F}_{\alpha}, g_{\beta \alpha})\) be two direct systems of sets relative to the same index set I; let \(\mathbf{E} = \varinjlim \mathbf{E}_{\alpha}\) , \(\mathbf{F} = \varinjlim \mathbf{F}_{\alpha}\) , and for each \(\alpha \in \mathbf{I}\) let \(f_{\alpha}\) (resp. \(g_{\alpha}\) ) be the canonical mapping of \(\mathbf{E}_{\alpha}\) (resp. \(\mathbf{F}_{\alpha}\) ) into \(\mathbf{E}\) (resp. \(\mathbf{F}\) ). For each \(\alpha \in \mathbf{I}\) let \(u_{\alpha}\) be a mapping of \(\mathbf{E}_{\alpha}\) into \(\mathbf{F}_{\alpha}\) such that, whenever \(\alpha \leqslant \beta\) , the diagram

\[
\begin{array}{c} \mathrm{E} _ {\alpha} \xrightarrow {u _ {\alpha}} \mathrm{F} _ {\alpha} \\ f _ {\beta \alpha} \Big \downarrow \qquad \qquad \qquad \qquad \qquad \qquad \qquad \qquad \qquad \Big \downarrow g _ {\beta \alpha} \\ \mathrm{E} _ {\beta} \xrightarrow [ u _ {\beta} ]{} \mathrm{F} _ {\beta} \end{array}
\]

is commutative. Then there exists a unique mapping \(u: \mathbf{E} \to \mathbf{F}\) such that, for each \(\alpha \in I\) , the diagram

\[
\begin{array}{c} \mathrm{E} _ {\alpha} \xrightarrow {u _ {\alpha}} \mathrm{F} _ {\alpha} \\ f _ {\alpha} \Big \downarrow \qquad \qquad \qquad \qquad \Big \downarrow g _ {\alpha} \\ \mathrm{E} \xrightarrow {} \mathrm{F} \\ u \end{array}
\]

is commutative.

Put \(v_{\alpha} = g_{\alpha} \circ u_{\alpha}\) . If \(\alpha \leqslant \beta\) , then by (22) we have

\[
v _ {\beta} \circ f _ {\beta \alpha} = g _ {\beta} \circ u _ {\beta} \circ f _ {\beta \alpha} = g _ {\beta} \circ g _ {\beta \alpha} \circ u _ {\alpha} = g _ {\alpha} \circ u _ {\alpha} = v _ {\alpha}.
\]

We may therefore apply Proposition 6 to the mappings \(v_{\alpha}\) , whence the existence and uniqueness of a mapping \(u: \mathbf{E} \to \mathbf{F}\) such that

\[
u \circ f _ {\alpha} = v _ {\alpha} = g _ {\alpha} \circ u _ {\alpha}
\]

for all \(\alpha \in I\) .

A family of mappings \(u_{\alpha} \colon \mathbf{E}_{\alpha} \to \mathbf{F}_{\alpha}\) which satisfies the conditions of Corollary 1 is called a direct system of mappings of \((\mathbf{E}_{\alpha}, f_{\beta \alpha})\) into \((\mathbf{F}_{\alpha}, g_{\beta \alpha})\) . The mapping defined in Corollary 1 is called the direct limit of the family \((u_{\alpha})\) and is written \(u = \lim_{\rightarrow} u_{\alpha}\) when there is no risk of ambiguity.

COROLLARY 2. Let \((\mathbf{E}_{\alpha}, f_{\beta \alpha})\) , \((\mathbf{F}_{\alpha}, g_{\beta \alpha})\) , \((\mathbf{G}_{\alpha}, h_{\beta \alpha})\) be three direct systems of sets relative to I. Let \(\mathbf{E} = \varinjlim \mathbf{E}_{\alpha}\) , \(\mathbf{F} = \varinjlim \mathbf{F}_{\alpha}\) , \(\mathbf{G} = \varinjlim \mathbf{G}_{\alpha}\) , and let \(f_{\alpha}\) (resp. \(g_{\alpha}, h_{\alpha}\) ) be the canonical mapping of \(\overrightarrow{\mathbf{E}_{\alpha}}\) (resp. \(\mathbf{F}_{\alpha}, \overrightarrow{\mathbf{G}_{\alpha}}\) ) into E (resp. F, G). If \((u_{\alpha})\) and \((v_{\alpha})\) are two direct systems of mappings \(u_{\alpha}: \mathbf{E}_{\alpha} \to \mathbf{F}_{\alpha}\) , \(v_{\alpha}: \mathbf{F}_{\alpha} \to \mathbf{G}_{\alpha}\) , then the mappings \(v_{\alpha} \circ u_{\alpha}: \mathbf{E}_{\alpha} \to \mathbf{G}_{\alpha}\) form a direct system of mappings, and we have

\[
\lim _ {\rightarrow} \left(v _ {\alpha} \circ u _ {\alpha}\right) = (\lim _ {\rightarrow} v _ {\alpha}) \circ (\lim _ {\rightarrow} u _ {\alpha}).\tag{25}
\]

For if we put \(w_{\alpha} = v_{\alpha} \circ u_{\alpha}\) , then for \(\alpha \leqslant \beta\) we have

\[
h _ {\beta \alpha} \circ w _ {\alpha} = \left(h _ {\beta \alpha} \circ v _ {\alpha}\right) \circ u _ {\alpha} = \left(v _ {\beta} \circ g _ {\beta \alpha}\right) \circ u _ {\alpha} = v _ {\beta} \circ \left(u _ {\beta} \circ f _ {\beta x}\right) = w _ {\beta} \circ f _ {\beta x},
\]

which shows that \((w_{\alpha})\) is a direct system of mappings. Furthermore, if \(u = \varinjlim u_{\alpha}\) and \(v = \varinjlim v_{\alpha}\) , then for all \(\alpha \in I\) we have

\[
(v \circ u) \circ f _ {\alpha} = v \circ (g _ {\alpha} \circ u _ {\alpha}) = h _ {\alpha} \circ (v _ {\alpha} \circ u _ {\alpha}),
\]

and by virtue of the uniqueness of the direct limit, we have \(v \circ u = \lim_{\rightarrow} w_{\alpha}\) .

PROPOSITION 7. Let \((\mathbf{E}_{\alpha}, f_{\beta \alpha})\) and \((\mathbf{E}_{\alpha}', f_{\beta \alpha}')\) be two direct systems of sets relative to I, and for each \(\alpha \in I\) let \(u_{\alpha}\) be a mapping of \(\mathbf{E}_{\alpha}\) into \(\mathbf{E}_{\alpha}'\) such that the \(u_{\alpha}\) form a direct system of mappings. Let \(u = \lim u_{\alpha}\) . If each \(u_{\alpha}\) is injective (resp. surjective) then \(u\) is injective (resp. surjective).

Let \(\mathbf{E} = \lim_{\rightarrow}\mathbf{E}_{\alpha},\mathbf{E}' = \lim_{\rightarrow}\mathbf{E}_{\alpha}'\) and let \(f_{\alpha}\colon \mathbf{E}_{\alpha}\to \mathbf{E},f_{\alpha}^{\prime}\colon \mathbf{E}_{\alpha}^{\prime}\to \mathbf{E}^{\prime}\) be the canonical mappings. Suppose that each \(u_{\alpha}\) is injective. To show that \(u\) is injective it is enough, by Proposition 6, to verify that if \(x\in \mathbf{E}_{\alpha}\) and \(y\in \mathbf{E}_{\alpha}\) are such that \(f_{\alpha}^{\prime}(u_{\alpha}(x)) = f_{\alpha}^{\prime}(u_{\alpha}(y))\) , then there exists \(\beta \geqslant \alpha\) such that \(f_{\beta \alpha}(x) = f_{\beta \alpha}(y)\) . Now the hypothesis implies (no. 6, Lemma 1) that there exists \(\beta \geqslant \alpha\) such that

\[
f _ {\beta \alpha} ^ {\prime} (u _ {\alpha} (x)) = f _ {\beta \alpha} ^ {\prime} (u _ {\alpha} (y)), \quad \text { i.e., } \quad u _ {\beta} (f _ {\beta \alpha} (x)) = u _ {\beta} (f _ {\beta \alpha} (y)),
\]

and hence \(f_{\beta \alpha}(x) = f_{\beta \alpha}(y)\) since \(u_{\beta}\) is injective.

Now suppose that the \(u_{\alpha}\) are surjective. Then we have

\[
\mathrm{E} ^ {\prime} = \bigcup_ {\alpha} f _ {\alpha} ^ {\prime} (\mathrm{E} _ {\alpha} ^ {\prime}) = \bigcup_ {\alpha} f _ {\alpha} ^ {\prime} (u _ {\alpha} (\mathrm{E} _ {\alpha})) = \bigcup_ {\alpha} u (f _ {\alpha} (\mathrm{E} _ {\alpha})) = u \left(\bigcup_ {\alpha} f _ {\alpha} (\mathrm{E} _ {\alpha})\right) = u (\mathrm{E}).
\]

With the notation of Proposition 7, let \(\mathbf{M}_{\alpha}\) be a subset of \(\mathbf{E}_{\alpha}\) for each \(\alpha \in \mathbf{I}\) ; if we have \(f_{\beta \alpha}(\mathbf{M}_{\alpha}) \in \mathbf{M}_{\beta}\) whenever \(\alpha \leqslant \beta\) , the family \((\mathbf{M}_{\alpha})_{\alpha \in \mathbf{I}}\) is said to be a direct system of subsets of the \(\mathbf{E}_{\alpha}\) . Let \(g_{\beta \alpha}\) (where \(\alpha \leqslant \beta\) ) be the mapping of \(\mathbf{M}_{\alpha}\) into \(\mathbf{M}_{\beta}\) whose graph is the same as that of the restriction of \(f_{\beta \alpha}\) to \(\mathbf{M}_{\alpha}\) . Then it is clear that \((\mathbf{M}_{\alpha}, g_{\beta \alpha})\) is a direct system of sets; and Proposition 7, applied to the canonical injections \(j_{\alpha}: \mathbf{M}_{\alpha} \to \mathbf{E}_{\alpha}\) , allows us to identify \(\mathbf{M} = \lim \mathbf{M}_{\alpha}\) with a subset of \(\mathbf{E}\) by means of the injection \(j = \lim j_{\alpha}\) .

COROLLARY. Let \((\mathbf{E}_{\alpha}, f_{\beta\alpha})\) and \((\mathbf{E}_{\alpha}^{\prime}, f_{\beta\alpha}^{\prime})\) be two direct systems of sets, let \((u_{\alpha})\) be a direct system of mappings \(u_{\alpha}: E_{\alpha} \to E_{\alpha}^{\prime}\) , and let \(u = \lim u_{\alpha}\) .

\begin{enumerate}
\def\labelenumi{(\roman{enumi})}
\tightlist
\item
  Let \((\mathbf{M}_{\alpha})\) be a direct system of subsets of the \(\mathbf{E}_{\alpha}\) . Then \((u_{\alpha}(\mathbf{M}_{\alpha}))\) is a direct system of subsets of the \(\mathbf{E}_{\alpha}'\) , and we have
\end{enumerate}

\[
\lim _ {\rightarrow} u _ {\alpha} (\mathbf {M} _ {\alpha}) = u (\lim _ {\rightarrow} \mathbf {M} _ {\alpha}).\tag{26}
\]

\begin{enumerate}
\def\labelenumi{(\roman{enumi})}
\setcounter{enumi}{1}
\tightlist
\item
  Let \((a_{\alpha}^{\prime})_{\alpha \in \mathbf{I}}\) be a family such that \(a_{\alpha}^{\prime} \in \mathrm{E}_{\alpha}^{\prime}\) for each \(\alpha \in \mathbf{I}\) and \(f_{\beta \alpha}^{\prime}(a_{\alpha}^{\prime}) = a_{\beta}^{\prime}\) whenever \(\alpha \leqslant \beta\) . Then the sets \(-u_{\alpha}^{1}(a_{\alpha}^{\prime})\) form a direct system of subsets of the \(\mathbf{E}_{\alpha}\) , and we have
\end{enumerate}

\[
\lim _ {\rightarrow} \overline {{{u}}} _ {\alpha} ^ {1} (a _ {\alpha} ^ {\prime}) = \overline {{{u}}} ^ {1} (a ^ {\prime}),\tag{27}
\]

where \(a'\) is the unique element of \(\lim_{\rightarrow} \mathbf{E}_{\alpha}'\) which is the canonical image of \(a_{\alpha}'\) for each \(\alpha \in \mathbf{I}\) .

\begin{enumerate}
\def\labelenumi{(\roman{enumi})}
\item
  It is evident that the \(u_{\alpha}(\mathbf{M}_{\alpha})\) form a direct system of subsets of the \(\mathbf{E}_{\alpha}'\) , and we may write \(u_{\alpha}(\mathbf{M}_{\alpha}) = v_{\alpha}(\mathbf{M}_{\alpha})\) , where \(v_{\alpha}\) is the mapping of \(\mathbf{M}_{\alpha}\) onto \(u_{\alpha}(\mathbf{M}_{\alpha})\) whose graph is the same as that of the restriction of \(u_{\alpha}\) to \(\mathbf{M}_{\alpha}\) . The formula (26) then follows from Proposition 7 because the \(v_{\alpha}\) are surjective.
\item
  Let \(\mathbf{N}_{\alpha} = \overline{u}_{\alpha}^{1}(a_{\alpha}^{\prime})\) . If \(\alpha \leqslant \beta\) and \(x_{\alpha} \in \mathbf{N}_{\alpha}\) , then
\end{enumerate}

\[
u _ {\beta} (f _ {\beta \alpha} (x _ {\alpha})) = f _ {\beta \alpha} ^ {\prime} (u _ {\alpha} (x _ {\alpha}) = f _ {\beta \alpha} ^ {\prime} (a _ {\alpha} ^ {\prime}) = a _ {\beta} ^ {\prime};
\]

hence \(f_{\beta \alpha}(x_{\alpha}) \in \mathbb{N}_{\beta}\) , and the \(\mathbb{N}_{\alpha}\) therefore form a direct system of subsets of the \(\mathbf{E}_{\alpha}\) . With the notation of the proof of Proposition 7, consider an element \(x \in \lim \mathbb{N}_{\alpha}\) . There exists \(\alpha \in \mathbb{I}\) and \(x_{\alpha} \in \mathbb{N}_{\alpha}\) such that \(x = f_{\alpha}(x_{\alpha})\) , so that \(u(x) = u(f_{\alpha}(x_{\alpha})) = f_{\alpha}'(u_{\alpha}(x_{\alpha})) = f_{\alpha}'(a_{\alpha}') = a'\) . Conversely, if \(x \in \overline{u}^1(a')\) and if \(x = f_\alpha(x_\alpha)\) for some \(\alpha \in I\) and some \(x_\alpha \in E_\alpha\) , then we have \(a' = u(f_\alpha(x_\alpha)) = f_\alpha'(u_\alpha(x_\alpha)) = f_\alpha'(a_\alpha')\) . Hence (no. 5, Lemma 1) there exists \(\beta \geqslant \alpha\) such that \(f_{\beta\alpha}'(u_\alpha'(x_\alpha)) = f_{\beta\alpha}'(a_\alpha') = a_\beta'\) ; i.e., \(u_\alpha(f_{\beta\alpha}(x_\alpha)) = a_\beta'\) , and therefore \(f_{\beta\alpha}(x_\alpha) \in N_\beta\) . Since \(x = f_\beta(f_{\beta\alpha}(x_\alpha))\) , it follows that \(x \in \lim_{\rightarrow} N_\alpha\) .

Remark. Suppose that, for each \(\alpha \in I\) , \(u_{\alpha}: E_{\alpha} \to E'\) is a mapping such that the family \((u_{\alpha})\) satisfies (23). Consider the direct system \((E_{\alpha}, i_{\beta \alpha})\) relative to I, where \(E_{\alpha}' = E'\) for all \(\alpha \in I\) , and \(i_{\beta \alpha}\) is the identity mapping of \(E'\) . Then (no. 5, Example 2) \(E'\) can be identified canonically with \(\lim E_{\alpha}'\) . If \(u_{\alpha}\) is considered as a mapping of \(E_{\alpha}\) into \(E_{\alpha}'\) , then \((u_{\alpha})\) is a direct system of mappings, and the mapping \(u: E \to E'\) defined by (24) is identified with the direct limit of this system of mappings. Hence, by abuse of language, we write \(u = \lim u_{\alpha}\) .

If J is a subset of I which is directed (with respect to the induced pre-ordering), it is clear that the pair consisting of the subfamily \((\mathrm{E}_{\alpha})_{\alpha\in\mathbb{J}}\) and the family \((f_{\beta\alpha})\) , where \(\alpha\leqslant\beta\) and \(\alpha\in J\) and \(\beta\in J\) , is a direct system of sets relative to J; it is said to be obtained by restricting the index set to J. Let E, \(E'\) respectively denote the direct limits of the families \((\mathrm{E}_{\alpha})_{\alpha\in\mathbb{I}}\) and \((\mathrm{E}_{\alpha})_{\alpha\in\mathbb{J}}\) , and for each \(\alpha\in I\) let \(f_{\alpha}:E_{\alpha}\to E\) denote the canonical mapping. Then \((f_{\alpha})_{\alpha\in\mathbb{J}}\) is a direct system of mappings, and consequently \(g=\lim f_{\alpha}\) is a mapping of \(E'\) into E, called canonical. Moreover, if \(J'\) is a directed subset of J, if \(E''\) is the direct limit of the family \((\mathrm{E}_{\alpha})_{\alpha\in\mathbb{J}'}\) , and if

\[
g ^ {\prime}: \quad \mathbf {E} ^ {\prime \prime} \rightarrow \mathbf {E} ^ {\prime} \quad \text { and } \quad g ^ {\prime \prime}: \mathbf {E} ^ {\prime \prime} \rightarrow \mathbf {E}
\]

are the canonical mappings, then it follows immediately from Proposition 6 that

\[
g ^ {\prime \prime} = g \circ g ^ {\prime}.\tag{28}
\]

PROPOSITION 8. Let I be a directed set, let \((\mathbf{E}_{\alpha}, f_{\beta \alpha})\) be a direct system of sets relative to I, and let \(\mathbf{E} = \lim \mathbf{E}_{\alpha}\) be its direct limit. Let J be a cofinal subset of I and let \(\mathbf{E}'\) be the direct limit of the direct system of sets obtained from \((\mathbf{E}_{\alpha}, f_{\beta \alpha})\) by restricting the index set to J. Then the canonical mapping \(g\) of \(\mathbf{E}'\) into E is bijective.

J is necessarily a directed set (§ 1, no. 10). We shall use the criteria of Proposition 6 to show that g is bijective. The condition for injectivity follows immediately from the definitions and from Lemma 1 of no. 5. To show that g is surjective, we note that for each \(\alpha \in J\) we have

\[
g (\mathbf {E} _ {\alpha}) = f _ {\alpha} (\mathbf {E} _ {\alpha}).
\]

Now, for each \(\beta \in I\) , there exists \(\gamma \in J\) such that \(\beta \leqslant \gamma\) , from which it follows that \(g(E_{\gamma}) \supset g(f_{\gamma \beta}(E_{\beta})) = f_{\beta}(E_{\beta})\) . E is therefore the union of the sets \(g(E_{\alpha})\) as \(\alpha\) runs through J.

